\subsubsection*{Perturbation arguments}
From the above limiting behaviours, we try to build a picture for $d= 1, 2$ and $3$. To construct the full scaling curve, AALR made the bold assumption that $\beta(g)$ is both continuous and monotonic. While mathematically this is an ansatz,but it is physically deeply motivated. Monotonicity reflects the fact that small conductance implies stronger localization and vice versa. AALR argued that the scaling function should be continuous because it describes how the conductance evolves when finite blocks of sites are combined. A cusp at finite conductance would introduce a non-smooth feature into this scale-evolution relation, for which they identified no physical justification. Crucially, this assumption directly contradicts Mott's hypothesis of a discontinuous conductivity jump\cite{AALR1979}.

In the ohmic regime, we never considered any quantum interference effects. Hence, to create a realistic scaling curve we introduce the leading order correction due to quantum interference of scattered waves. Lets begin with the free electron gas with few disorder and treating those disorder as perturbation, using linear response theory to an applied m=electromagnetic field, we can get the conductance expressed in terms of current-current correlation with the help of Kubo formalism (see Appendix 1).
\begin{equation}
    \sigma_{\alpha\beta}
    (\mathbf r,\mathbf r',\omega)
    =
    \frac{i e^2}{\omega}
    \chi_{\alpha\beta}^{R}
    (\mathbf r,\mathbf r',\omega)
    +
    \frac{e^2 n(\mathbf r)}{m\omega}
    \delta(\mathbf r-\mathbf r')
    \delta_{\alpha\beta},
    \label{eq:kubo_conductivity}
\end{equation}
where $n(\mathbf r)$ is the particle density, $m$ is the electron
mass. The function
$\chi_{\alpha\beta}^{R}$ denotes the retarded correlation function
of the paramagnet current density. In the time domain, this retarded current--current correlation
function is defined as
\begin{equation}
    \chi_{\alpha\beta}^{R}
    (\mathbf r,\mathbf r',t-t')
    =
    -i\theta(t-t')
    \left\langle
    \left[
        \hat J_0^\alpha(\mathbf r,t),
        \hat J_0^\beta(\mathbf r',t')
    \right]
    \right\rangle_0,
    \label{eq:retarded_current_correlation}
\end{equation}
where $\theta(t-t')$ is the Heaviside step function,
$\hat J_0^\alpha$ is the $\alpha$ component of the paramagnetic current-density operator, the subscript "$0$" denotes expectation value is taken with respect tot the unperturbed hamiltonian.\\
The correlation can be demonstrated as the 
Langer and Neal later in 1966, gave theory to higher order term in impurity density for the disorder perturbation. They used diagrammatic approach to the problem and gave the amplitude of the process which had a singularity at 
 